% answers_short.tex —— 1987 年·全国卷·文科 · 简短答案速查
% 由 tools/make_short_answers.py 自 answers.tex 的 \gksolhead/\gkanswer 提取，
% 并与 1987/1987_数学真题_全国文_答案要点.md 逐题比对后生成，勿手改（改源 MD 后重跑）。

\begin{gknotice}
\textbf{使用说明}：本节只给最终答案便于快速核对，完整推导见解析版。答案为编写者独立推导，\textbf{非官方标准答案}。
\end{gknotice}

\gkpart{题目}
\begin{itemize}[leftmargin=2.6em]
  \item[\textbf{1.}] D
  \item[\textbf{2.}] C
  \item[\textbf{3.}] B
  \item[\textbf{4.}] D
  \item[\textbf{5.}] C
  \item[\textbf{6.}] B
  \item[\textbf{7.}] A
  \item[\textbf{8.}] D
  \item[\textbf{9.}] $\dfrac\pi2$
  \item[\textbf{10.}] $\lambda\in(-\infty,-2)\cup(-1,+\infty)$
  \item[\textbf{11.}] $n=8$
  \item[\textbf{12.}] $2$
  \item[\textbf{13.}] $72$
  \item[\textbf{14.}] $\left(-\dfrac13,1\right)$
  \item[\textbf{15.}] $3\pi$
  \item[\textbf{16.}] $0$
  \item[\textbf{17.}] $\displaystyle \frac12-\frac{\sqrt3}{2}i$ 或 $\displaystyle 2+\sqrt3\,i$
  \item[\textbf{18.}] 见解析正文。
  \item[\textbf{19.}] $0<a<1$
  \item[\textbf{20.}] (1) 见解析正文；(2) 见解析正文
  \item[\textbf{21.}] (1) $\displaystyle a_n=\frac{(-k)^{n-1}}{(1-k)^n}$；(2) $k<\dfrac12$
  \item[\textbf{22.}] $18$ 或 $50$
\end{itemize}
